\section*{AI4PowerGrids domain checklist}
% 中文：
%   AI4PowerGrids 电网领域检查表

\begin{enumerate}
  \item \textbf{Physics feasibility.}
  \emph{Not applicable to AC power-flow feasibility.} CityLearn enforces the
  simulated battery, HVAC, and building thermal dynamics and bounds the control
  actions, while Section~4 and Appendices~B--C report thermal-comfort violations
  explicitly. The benchmark does not include a distribution-network model;
  consequently, voltage, line-flow, and full AC power-flow feasibility are not
  evaluated and no such claim is made.
  % 中文（物理可行性）：
  %   不适用于交流潮流可行性。CityLearn 会执行模拟电池、HVAC 和建筑热动态，
  %   并限制控制动作；第 4 节及附录 B--C 明确报告了室内热舒适度违规情况。
  %   该基准不包含配电网络模型，因此本文没有评估电压、线路潮流或完整交流潮流
  %   可行性，也不作出相关声明。

  \item \textbf{Out-of-distribution evaluation.}
  Section~5 and Appendix~D evaluate the same feature-design principle in two
  seasonal/climate regimes (Vermont heating and Texas cooling), each with a
  chronological held-out month. The Texas controller is retrained, so this is
  not zero-shot evaluation of one policy on an unseen climate or topology. This
  missing test is stated as a limitation.
  % 中文（分布外评估）：
  %   第 5 节和附录 D 在 Vermont 供暖季与 Texas 制冷季两种气候工况中评估
  %   同一套特征设计原则，并分别使用按时间顺序留出的测试月份。
  %   Texas 控制器经过重新训练，因此这不是将同一策略零样本部署到未见气候
  %   或拓扑上的评估；本文已将缺少这种测试明确列为局限。

  \item \textbf{Failure modes.}
  Section~4 reports the reversal between mean and worst-building comfort and
  the unequal numbers of training episodes. Appendices~B--C expose daily metric variation
  and persistent temperature violations in individual homes. Section~3 also
  explains why joint router/expert training was rejected after routing
  concentration and policy instability were observed.
  % 中文（失效模式）：
  %   第 4 节报告了平均舒适度与最差建筑舒适度之间的结论反转，以及各方法训练
  %   episode 数不同的问题。附录 B--C 展示逐日指标波动和个别住宅持续温度超限。
  %   第 3 节还说明：由于观察到路由集中和策略不稳定，最终没有采用路由器与
  %   专家联合训练方案。

  \item \textbf{System-level effect.}
  Section~4 and Appendix~A report portfolio CV-RMSE, NMBE, peak demand, and
  ramping together with building-level comfort. This evaluates each controller
  inside the complete closed-loop CityLearn portfolio rather than as a
  standalone predictor or router.
  % 中文（系统级影响）：
  %   第 4 节和附录 A 同时报告组合层面的 CV-RMSE、NMBE、峰值需求、ramping，
  %   以及建筑层面的舒适度。由此评估的是完整 CityLearn 建筑组合闭环中的控制器，
  %   而不是脱离系统的独立预测器或路由器。

  \item \textbf{Data and code.}
  The experiments use the open CityLearn implementation and IEA EBC Annex~96
  Common Exercise~1 benchmark data. We plan to provide the training and
  evaluation code, environment specifications, configuration scripts, processed
  metric tables, and plotting code as anonymized supplementary material during
  review, followed by a de-anonymized public repository upon acceptance.
  % 中文（数据与代码）：
  %   实验使用开放的 CityLearn 实现和 IEA EBC Annex 96 Common Exercise 1
  %   基准数据。我们计划在评审期间以匿名补充材料形式提供训练与评估代码、环境
  %   规格、配置脚本、处理后的指标表和绘图代码；若论文录用，再公开去匿名化后的
  %   代码仓库。
\end{enumerate}
